\documentclass[aspectratio=169,10pt]{beamer}

% Компилировать XeLaTeX или LuaLaTeX (fontspec + polyglossia)
\usepackage{fontspec}
\usepackage{polyglossia}
\setmainlanguage{russian}
\setotherlanguage{english}
\setmainfont{DejaVu Sans}[ItalicFont={DejaVu Sans},ItalicFeatures={FakeSlant=0.22},BoldItalicFont={DejaVu Sans},BoldItalicFeatures={FakeSlant=0.22}]
\setsansfont{DejaVu Sans}[ItalicFont={DejaVu Sans},ItalicFeatures={FakeSlant=0.22},BoldItalicFont={DejaVu Sans},BoldItalicFeatures={FakeSlant=0.22}]
\setmonofont{DejaVu Sans Mono}[Scale=0.88]
\newfontfamily\cyrillicfont{DejaVu Sans}[ItalicFont={DejaVu Sans},ItalicFeatures={FakeSlant=0.22},BoldItalicFont={DejaVu Sans},BoldItalicFeatures={FakeSlant=0.22}]
\newfontfamily\cyrillicfontsf{DejaVu Sans}[ItalicFont={DejaVu Sans},ItalicFeatures={FakeSlant=0.22},BoldItalicFont={DejaVu Sans},BoldItalicFeatures={FakeSlant=0.22}]
\newfontfamily\cyrillicfonttt{DejaVu Sans Mono}[Scale=0.88]

\usepackage{graphicx}
\usepackage{tikz}
\usetikzlibrary{arrows.meta,shapes.multipart,calc,positioning}

\definecolor{main}{HTML}{1F4E79}
\definecolor{accent}{HTML}{B5472B}
\setbeamercolor{structure}{fg=main}
\setbeamercolor{frametitle}{fg=white,bg=main}
\setbeamercolor{block title}{fg=white,bg=main}
\setbeamercolor{block body}{bg=main!7}
\setbeamercolor{block title alerted}{fg=white,bg=accent}
\setbeamercolor{block body alerted}{bg=accent!8}
\setbeamerfont{frametitle}{size=\large,series=\bfseries}
\setbeamerfont{block title}{size=\small,series=\bfseries}
\setbeamerfont{block body}{size=\small}
\setbeamerfont{block body alerted}{size=\small}
\setbeamerfont{itemize/enumerate body}{size=\small}
\setbeamerfont{itemize/enumerate subbody}{size=\footnotesize}
\setbeamertemplate{navigation symbols}{}
\setbeamertemplate{items}[circle]
\setbeamersize{text margin left=0.8cm,text margin right=0.8cm}
\setbeamertemplate{frametitle}{%
  \nointerlineskip
  \begin{beamercolorbox}[wd=\paperwidth,ht=1.05cm,dp=0.35cm,leftskip=0.8cm,rightskip=0.8cm]{frametitle}
    \usebeamerfont{frametitle}\insertframetitle
  \end{beamercolorbox}%
  \vskip1.5mm
}
\setbeamertemplate{footline}{%
  \hbox to\paperwidth{\hspace{0.8cm}\color{black!55}\scriptsize
    Порождающие паттерны проектирования\hfill
    \insertframenumber\,/\,\inserttotalframenumber\hspace{0.8cm}}%
  \vskip3mm
}

\tikzset{
  uml/.style={draw=black!70, fill=main!6, rectangle split,
    rectangle split draw splits, font=\footnotesize,
    text width=3.2cm, inner sep=3pt},
  cls3/.style={uml, rectangle split parts=3,
    rectangle split part align={center,left,left}},
  cls2/.style={uml, rectangle split parts=2,
    rectangle split part align={center,left}},
  cls1/.style={uml, rectangle split parts=1, align=center},
  % уменьшенный шрифт: применяется ПОСЛЕ cls*, без глобальных побочных эффектов
  scls1/.style={cls1, font=\scriptsize},
  scls2/.style={cls2, font=\scriptsize},
  scls3/.style={cls3, font=\scriptsize},
  % обобщение: сплошная линия, незакрашенный треугольник
  inherit/.style={-{Triangle[open,length=9pt,width=8pt]}, thick},
  % реализация: пунктир, незакрашенный треугольник
  realize/.style={-{Triangle[open,length=9pt,width=8pt]}, thick, dashed},
  % навигируемая ассоциация: сплошная линия, открытая «галочка» (V)
  assoc/.style={thick, -{Straight Barb[length=8pt,width=6pt]}},
  % зависимость: пунктир, открытая «галочка» (V)
  depend/.style={-{Straight Barb[length=9pt,width=7pt]}, thick, dashed},
  % композиция: закрашенный ромб у целого
  compo/.style={{Diamond[length=10pt,width=7pt]}-, thick},
  ndash/.style={dashed, black!60},
  lbl/.style={font=\scriptsize, inner sep=1.5pt, fill=white, text=black!80, align=center},
  note/.style={draw=black!50, fill=yellow!22, font=\scriptsize\ttfamily,
    align=left, inner sep=3pt},
}

\title{Порождающие паттерны проектирования}
\date{\today}

\begin{document}

\begin{frame}[plain]
  \vfill
  \centering
  {\color{main}\fontsize{26}{31}\selectfont\bfseries Порождающие паттерны\\проектирования\par}
  \vspace{5mm}
  {\color{accent}\rule{6cm}{1.2pt}\par}
  \vspace{5mm}
  {\large UML-диаграммы\par}
  \vspace{7mm}
  {\small Singleton \,·\, Abstract Factory \,·\, Builder \,·\, Factory Method \,·\,
   Lazy initialization\\[1mm]
   Multiton \,·\, Object pool \,·\, Prototype\par}
  \vfill
\end{frame}

\begin{frame}{Содержание}
  \begin{columns}[T]
    \begin{column}{0.46\textwidth}
      \begin{enumerate}\setlength\itemsep{3pt}
        \item Singleton
        \item Abstract Factory
        \item Builder
        \item Factory Method
        \item Lazy initialization
      \end{enumerate}
    \end{column}
    \begin{column}{0.46\textwidth}
      \begin{enumerate}\setcounter{enumi}{5}\setlength\itemsep{3pt}
        \item Multiton
        \item Object pool
        \item Prototype
      \end{enumerate}
    \end{column}
  \end{columns}
\end{frame}

% --------------------------------------------------------------- Слайд 3
\begin{frame}{Singleton: UML-диаграмма}
  \centering
  \begin{tikzpicture}[scale=1.1, transform shape]
    \node[cls1, text width=1.4cm] (client) at (0,0.8) {\textbf{Client}};
    \node[cls3, text width=4.0cm] (s) at (6.2,0.8) {\textbf{Singleton}
      \nodepart{two} − \underline{instance : Singleton}
      \nodepart{three} − Singleton()\\
      + \underline{getInstance() : Singleton}};
    \draw[depend] (client.east) -- (s.west)
      node[lbl, midway, above=1mm] {getInstance()};
    \node[note] (n) at (6.2,-1.3) {if (instance == null)\\
      \quad instance = new Singleton();\\
      return instance;};
    \draw[ndash] (n.north) -- (s.south);
  \end{tikzpicture}

  \vspace{2mm}
  {\small Доступ только через статический \texttt{getInstance()}; конструктор закрыт.
   Экземпляр создаётся лениво и существует в единственном числе.}
\end{frame}

% --------------------------------------------------------------- Слайд 4
\begin{frame}{Abstract Factory: UML-диаграмма}
  \centering
  \resizebox{\textwidth}{!}{%
  \begin{tikzpicture}
    \node[scls1, text width=1.2cm] (client) at (-0.3,4.4) {\textbf{Client}};
    \node[scls2, text width=3.0cm] (af) at (3.1,4.4)
      {«interface»\\\textbf{AbstractFactory}
       \nodepart{two}+ createProductA() : AbstractProductA\\+ createProductB() : AbstractProductB};
    \node[scls2, text width=2.8cm] (cf1) at (1.6,1.4)
      {\textbf{ConcreteFactory1}
       \nodepart{two}+ createProductA() : ProductA1\\+ createProductB() : ProductB1};
    \node[scls2, text width=2.8cm] (cf2) at (4.7,1.4)
      {\textbf{ConcreteFactory2}
       \nodepart{two}+ createProductA() : ProductA2\\+ createProductB() : ProductB2};
    \node[scls2, text width=2.7cm] (apa) at (8.2,4.4)
      {«interface»\\\textbf{AbstractProductA}\nodepart{two}+ operationA()};
    \node[scls2, text width=2.7cm] (apb) at (11.4,4.4)
      {«interface»\\\textbf{AbstractProductB}\nodepart{two}+ operationB()};
    \node[scls1, text width=1.5cm] (a1) at (7.4,1.4) {\textbf{ProductA1}};
    \node[scls1, text width=1.5cm] (a2) at (9.2,1.4) {\textbf{ProductA2}};
    \node[scls1, text width=1.5cm] (b1) at (11.0,1.4) {\textbf{ProductB1}};
    \node[scls1, text width=1.5cm] (b2) at (12.8,1.4) {\textbf{ProductB2}};
    \draw[assoc] (client.east) -- (af.west);
    \draw[realize] ([xshift=-0.65cm]cf1.north) -- ([xshift=-0.55cm]af.south);
    \draw[realize] ([xshift=0.65cm]cf2.north) -- ([xshift=0.55cm]af.south);
    \draw[realize] ([xshift=-0.35cm]a1.north) -- ([xshift=-0.45cm]apa.south);
    \draw[realize] ([xshift=0.35cm]a2.north) -- ([xshift=0.45cm]apa.south);
    \draw[realize] ([xshift=-0.35cm]b1.north) -- ([xshift=-0.45cm]apb.south);
    \draw[realize] ([xshift=0.35cm]b2.north) -- ([xshift=0.45cm]apb.south);
    \draw[depend] (client.north) -- ++(0,1.0) -| (apa.north);
    \draw[depend] ([xshift=0.2cm]client.north) -- ++(0,1.35) -| (apb.north);
  \end{tikzpicture}}

  \vspace{2mm}
  {\small Клиент зависит от абстрактных типов; конкретная фабрика возвращает
   согласованную пару продуктов, ковариантно уточняя типы результатов методов.}
\end{frame}

% --------------------------------------------------------------- Слайд 5
\begin{frame}{Builder: UML-диаграмма}
  \centering
  \begin{tikzpicture}
    \node[cls3, text width=2.9cm] (dir) at (0,3.4) {\textbf{Director}
      \nodepart{two}− builder : Builder
      \nodepart{three}+ construct()};
    \node[cls2, text width=3.7cm] (b) at (5.8,3.4) {«interface»\\\textbf{Builder}
      \nodepart{two}+ buildPartA()\\+ buildPartB()\\+ getResult() : Product};
    \node[cls3, text width=3.7cm] (cb) at (5.8,0.4) {\textbf{ConcreteBuilder}
      \nodepart{two}− product : Product
      \nodepart{three}+ buildPartA()\\+ buildPartB()\\+ getResult() : Product};
    \node[cls2, text width=2.2cm] (p) at (11.2,0.4) {\textbf{Product}
      \nodepart{two}− parts : List};
    \draw[assoc] (dir.east) -- (b.west)
      node[lbl, midway, above=1mm] {builder};
    \draw[realize] (cb.north) -- (b.south);
    \draw[assoc] (cb.east) -- (p.west)
      node[lbl, midway, above=1mm] {product};
  \end{tikzpicture}

  \vspace{3mm}
  {\small Director задаёт порядок шагов, ConcreteBuilder собирает Product;
   готовый результат запрашивают у строителя.}
\end{frame}

% --------------------------------------------------------------- Слайд 6
\begin{frame}{Factory Method: UML-диаграмма}
  \centering
  \begin{tikzpicture}[scale=1.25, transform shape]
    \node[cls2, text width=4.0cm] (cr) at (2.4,3.2) {\textbf{\textit{Creator}}
      \nodepart{two}+ someOperation()\\+ \textit{factoryMethod() : Product}};
    \node[cls2, text width=4.0cm] (cc) at (2.4,0.5) {\textbf{ConcreteCreator}
      \nodepart{two}+ factoryMethod() : Product};
    \node[cls2, text width=2.75cm] (pr) at (9.6,3.2) {«interface»\\\textbf{Product}
      \nodepart{two}+ operation()};
    \node[cls2, text width=2.75cm] (cp) at (9.6,0.5) {\textbf{ConcreteProduct}
      \nodepart{two}+ operation()};
    \draw[inherit] (cc.north) -- (cr.south);
    \draw[realize] (cp.north) -- (pr.south);
    \draw[depend] (cr.east) -- (pr.west)
      node[lbl, midway, above=1mm] {«use»};
    \draw[depend] (cc.east) -- (cp.west)
      node[lbl, midway, above=1mm] {«create»};
  \end{tikzpicture}

  \vspace{3mm}
  {\small Creator \textbf{объявляет} абстрактный фабричный метод, возвращающий Product.
   ConcreteCreator реализует его и создаёт ConcreteProduct — подкласс сам решает,
   какой класс инстанцировать.}
\end{frame}

% --------------------------------------------------------------- Слайд 7
\begin{frame}{Lazy initialization: UML-диаграмма}
  \centering
  \begin{tikzpicture}
    \node[cls1, text width=1.4cm] (client) at (-1.1,1.6) {\textbf{Client}};
    \node[cls3, text width=4.7cm] (lz) at (4.6,1.6) {\textbf{LazyHolder}
      \nodepart{two}− resource : Resource = null
      \nodepart{three}+ getResource() : Resource\\− createResource() : Resource};
    \node[cls2, text width=2.4cm] (r) at (9.9,1.6) {\textbf{Resource}
      \nodepart{two}+ use()};
    \draw[depend] (client.east) -- (lz.west)
      node[lbl, midway, above=1mm] {getResource()};
    \draw[assoc] (lz.east) -- (r.west)
      node[lbl, pos=0.5, above=1mm] {resource};
    \node[note] (n) at (4.6,-0.9) {if (resource == null)\\
      \quad resource = createResource();\\
      return resource;};
    \draw[ndash] (n.north) -- (lz.south);
  \end{tikzpicture}

  \vspace{3mm}
  {\small Объект создаётся при первом обращении и затем переиспользуется;
   поле \texttt{resource} может не содержать ссылку до первого вызова.}
\end{frame}

% --------------------------------------------------------------- Слайд 8
\begin{frame}{Multiton: UML-диаграмма}
  \centering
  \begin{tikzpicture}[scale=1.0, transform shape]
    \node[cls1, text width=1.4cm] (client) at (-2.0,1.0) {\textbf{Client}};
    \node[cls3, text width=5.3cm] (m) at (4.8,1.0) {\textbf{Multiton}
      \nodepart{two}− \underline{instances : Map<Key, Multiton>}
      \nodepart{three}− Multiton(key : Key)\\
      + \underline{getInstance(key : Key) : Multiton}};
    \draw[depend] (client.east) -- (m.west)
      node[lbl, midway, above=1mm] {getInstance(key)};
    % Самоассоциация: целое одно (ромб), экземпляров 0..* — по одному на ключ
    \draw[compo] ([yshift=0.35cm]m.east) -- ++(0.8,0) -- ++(0,-1.0)
                 -- ([yshift=-0.65cm]m.east);
    \node[lbl, anchor=south west] at ([xshift=0.5cm,yshift=0.45cm]m.east) {1};
    \node[lbl, anchor=west] at (8.45,0.5) {instances};
    \node[lbl, anchor=north] at (7.95,0.28) {0..*};
    \node[note] (n) at (4.8,-1.15) {if (!instances.containsKey(key))\\
      \quad instances[key] = new Multiton(key);\\
      return instances[key];};
    \draw[ndash] (n.north) -- (m.south);
  \end{tikzpicture}

  \vspace{2mm}
  {\small По одному экземпляру на каждый ключ: повторный вызов \texttt{getInstance(key)}
   возвращает уже созданный объект. Всего экземпляров 0..*.}
\end{frame}

% --------------------------------------------------------------- Слайд 9
\begin{frame}{Object pool: UML-диаграмма}
  \centering
  \begin{tikzpicture}
    \node[cls1, text width=1.4cm] (client) at (-1.0,1.8) {\textbf{Client}};
    \node[cls3, text width=4.4cm] (pool) at (5.4,1.8) {\textbf{ObjectPool}
      \nodepart{two}− available : List<Reusable>\\− inUse : List<Reusable>
      \nodepart{three}+ acquire() : Reusable\\+ release(obj : Reusable)};
    \node[cls3, text width=2.6cm] (re) at (10.9,1.8) {\textbf{Reusable}
      \nodepart{two}− state : State
      \nodepart{three}+ reset()\\+ doWork()};
    \draw[depend] (client.east) -- (pool.west)
      node[lbl, midway, above=1mm] {acquire() / release()};
    \draw[compo] (pool.east) -- (re.west);
    \node[lbl, anchor=south west] at ([xshift=0.5cm,yshift=0.1cm]pool.east) {1};
    \node[lbl, anchor=north east] at ([xshift=-0.45cm,yshift=-0.1cm]re.west) {0..*};
    \draw[depend] (client.north) -- ++(0,1.5) -| (re.north)
      node[lbl, pos=0.25, above=1mm] {uses};
  \end{tikzpicture}

  \vspace{3mm}
  {\small Клиент берёт объект методом \texttt{acquire()} и возвращает через
   \texttt{release()}; пул владеет жизненным циклом объектов и сбрасывает
   их состояние вызовом \texttt{reset()}.}
\end{frame}

% --------------------------------------------------------------- Слайд 10
\begin{frame}{Prototype: UML-диаграмма}
  \centering
  \begin{tikzpicture}
    \node[cls3, text width=3.5cm] (cl) at (1.8,3.2) {\textbf{Client}
      \nodepart{two}− prototype : Prototype
      \nodepart{three}+ operation()};
    \node[cls2, text width=3.0cm] (pr) at (7.4,3.2) {«interface»\\\textbf{Prototype}
      \nodepart{two}+ clone() : Prototype};
    \node[cls3, text width=3.2cm] (p1) at (5.0,0.5) {\textbf{ConcretePrototype1}
      \nodepart{two}− field : Type
      \nodepart{three}+ clone() : Prototype};
    \node[cls3, text width=3.2cm] (p2) at (10.4,0.5) {\textbf{ConcretePrototype2}
      \nodepart{two}− field : Type
      \nodepart{three}+ clone() : Prototype};
    \node[note] (n) at (0.4,0.5) {return new\\ConcretePrototype1(this);};
    \draw[assoc] (cl.east) -- (pr.west)
      node[lbl, midway, above=1mm] {prototype};
    \draw[realize] (p1.north) -- ([xshift=-0.45cm]pr.south);
    \draw[realize] (p2.north) -- ([xshift=0.45cm]pr.south);
    \draw[ndash] (n.east) -- (p1.west);
  \end{tikzpicture}

  \vspace{3mm}
  {\small Клиент хранит ссылку на образец и клонирует его через интерфейс
   \texttt{Prototype}, не зная конкретного класса; поверхностное или глубокое
   копирование определяется реализацией \texttt{clone()}.}
\end{frame}

\end{document}
